% Pursuing Stacks (cotes 134-2 to 134-8): archive cote 134-6, site page 43 (PDF page 44), typescript page 361.
% typescript. Draft compared: ps5.tex:3485-3544.
% AI-transcribed by Claude Opus 5.5 (claude-opus-5-5, adaptive thinking, effort high), Message Batches API, batch ps3 (msgbatch_01XdFmuVVBAcgc7EFxdtCqsj), 2026-11-12; not yet checked by eye.
\scanpage{134-6/43}
\mspage{361}

equivalently,  to describe it's cohomology invariants with \xout{coefficients}
\xout{in} arbitrary \uline{constant}  coefficients.  Now among the strongest reflexes
I acquired in the past while working with cohomology,  was systematical
-ly to look at coefficients which are arbitrary sheaves  (abelian say),
and to view \deleted{at} constant or locally constant coefficients as being just
particular cases.  This reflex has been remaining id\xout{\illegible{1}}le,  not to say
repressed,  during nearly all of the reflections of the last four
months,  due to the fact that in the whole modelizing story woven \xout{around}
around weak equivalence, there was a rather exclusive emphasis on con-%
stant and locally constant coefficients,  disregarding any other coeffi%
cients throughout.  Probably,  while reflecting on abelianization,
a more or less underground reminiscence must have been around of the
semisimplicial boundary operators hav ing a meaning for computing coho-%
mology of``something''\hand{,} with coefficients in arbitrary sheaves  -  and
also that to get it straight,  one had to be careful not to get \xout{mixed}
mixed up in the variances.  B\over{\unsure{u}}{u}t I just did'nt want to dive into all \xout{thi}
this again if I could help \hand{\unsure{''}} -  and now it is getting clear,  after a
day or two of feeling silly,  that it can't be helped,  and I'll have
to write things down at last,  however ``well known'' they may be.

Let A be a small category.  In section 93 we defined an ``abelia-%
nizator'' for A \xout{for} to be a chain complex  $L$\over{,}{$_\bullet$}  in the additive envelope
Add($A^{o}$)  of the category $A^{o}$ opposite to A,  \xout{endowed with} satisfying a
suitable condition of commutativity (in the diagram (20) of p.328),
and endowed with a mild extra structure\inserted{$\lambda$,} expressing this commutativity.
The function of an abelianizator in loc. cit.  was essentially to allow
for a simultaneous ``computational'' description of the homology structu-%
re of the homotopy types stemming from a variable object X in $A^{\wedge}$,  or
equivalently,  to describe cohomology of such X (as an object of
a suitable derived category say, to get it at strongest) with coeffi-%
cients in any (\uline{constant}) \xout{group or ring} ring or abelian group.  Intro-%
ducing by an independant symbol the opposite category
\[ (1) \qquad B = A^{o} \qquad , \]
I want now to establish a relationship between this property or functi-%
on of a chain complex  $L.$\hand{\unsure{$_\bullet$}}  in Add(B),  \deleted{\xout{and}a} involving objects in $A^{\wedge}$
\margin{$A^{\wedge}_{ab}$ (circled)}and their abelianizations\inserted{in $A^{\wedge}_{ab}$}, with an apparently different one,  in
terms of a variable object of $B^{\wedge}_{ab}$  (\uline{not}  $A^{\wedge}_{ab}$  this time\xout{s} ! ),  namely
expressing \uline{cohomology of B} (i.e. of the topos $B^{\wedge}$ defined by B) with
coefficients in an \uline{arbitrary}\inserted{abelian} presheaf  F on B  i.e. an arbitrary object
\margin{43}
